% Simple penetration test report template. Compile with pdfLaTeX twice (or latexmk -pdf).
% Overleaf: upload this file and select pdfLaTeX as the compiler.
% Replace all bracketed prompts before submitting. No findings are pre-populated.
\documentclass[11pt,letterpaper]{article}

\usepackage[margin=1in]{geometry}
\usepackage[T1]{fontenc}
\usepackage{lmodern}
\usepackage{microtype}
\usepackage[table]{xcolor}
\usepackage{graphicx}
\usepackage{array,tabularx,longtable,booktabs}
\usepackage{enumitem}
\usepackage{listings}
\usepackage{fancyhdr}
\usepackage{hyperref}

% EDIT THESE FIRST. Use \& for an ampersand in names.
\newcommand{\ClientName}{[Client name]}
\newcommand{\TeamName}{[Team name]}
\newcommand{\EngagementName}{Penetration Test}
\newcommand{\TestingDates}{[Start date -- end date]}
\newcommand{\ReportDate}{[Report date]}
\newcommand{\ReportVersion}{0.1 -- Draft}

% Formatting: standard fonts, generous margins, restrained color.
\renewcommand{\familydefault}{\sfdefault}
\definecolor{accent}{HTML}{234A61}
\definecolor{muted}{HTML}{555555}
\definecolor{lightgray}{HTML}{F3F4F5}
\definecolor{critical}{HTML}{85213B}
\definecolor{high}{HTML}{A33125}
\definecolor{medium}{HTML}{805500}
\definecolor{low}{HTML}{245E85}
\definecolor{informational}{HTML}{555555}
\hypersetup{colorlinks=true,linkcolor=accent,urlcolor=accent,
  pdftitle={Penetration Test Report}}
\setlength{\parindent}{0pt}
\setlength{\parskip}{6pt}
\setlength{\headheight}{15pt}
\setlength{\emergencystretch}{2em}
\renewcommand{\arraystretch}{1.25}
\setlist{nosep,leftmargin=*}
\setcounter{tocdepth}{2}
\pagestyle{fancy}
\fancyhf{}
\fancyhead[L]{\small\textcolor{muted}{\ClientName}}
\fancyhead[R]{\small\textcolor{muted}{Penetration Test Report}}
\fancyfoot[L]{\small Confidential}
\fancyfoot[R]{\small\thepage}
\renewcommand{\headrulewidth}{0.3pt}
\newcolumntype{Y}{>{\raggedright\arraybackslash}X}
\newcommand{\prompt}[1]{\textit{\textcolor{muted}{[#1]}}}
\newcommand{\severity}[1]{\textbf{\textcolor{#1}{\MakeUppercase{#1}}}}
\lstset{basicstyle=\small\ttfamily,backgroundcolor=\color{lightgray},
  breaklines=true,columns=fullflexible,keepspaces=true,
  showstringspaces=false,frame=single,rulecolor=\color{lightgray},
  xleftmargin=6pt,xrightmargin=6pt,aboveskip=10pt,belowskip=10pt}

% Usage: \finding{F-01}{Title}{high}{finding:unique-label}
% Severity: critical, high, medium, low, or informational.
% Each finding starts on a new page; its contents can span multiple pages.
\newcommand{\finding}[4]{%
  \clearpage
  \subsection{#1: #2}\label{#4}
  \severity{#3}\par
}

\begin{document}

% COVER -----------------------------------------------------------------------
\begin{titlepage}
  {\large\TeamName\par}
  \vspace{1.6in}
  {\Huge\bfseries\color{accent}Penetration Test\\[6pt]Report\par}
  \vspace{0.4in}
  {\LARGE\ClientName\par}
  {\large\EngagementName\par}
  \vspace{0.6in}
  \begin{tabular}{@{}ll@{}}
    Testing period & \TestingDates \\
    Report date & \ReportDate \\
    Version & \ReportVersion \\
    Prepared by & \TeamName
  \end{tabular}
  \vfill
  \textbf{Confidential}\par
  Prepared for the client and authorized competition evaluators.
  Findings describe the environment observed during the testing period.
\end{titlepage}

\pagenumbering{roman}
\tableofcontents
\clearpage
\pagenumbering{arabic}

% EXECUTIVE SUMMARY -----------------------------------------------------------
\section{Executive Summary}
\prompt{In two or three sentences, explain what the client does, the objective
of the engagement, and which business services were assessed.}

\subsection*{Overall assessment}
\prompt{Explain the most important demonstrated risk in plain language.
State what access the team obtained and what business operations or data could
be affected. Separate observed results from potential consequences.}

\subsection*{Key observations}
\begin{itemize}
  \item \prompt{Most significant validated weakness and business impact.}
  \item \prompt{Another important weakness, or a chain of findings.}
  \item \prompt{A security control that worked, supported by testing.}
\end{itemize}

\subsection*{Recommended priorities}
\begin{enumerate}
  \item \prompt{Immediate action to reduce the greatest demonstrated risk.}
  \item \prompt{Near-term correction that addresses a shared root cause.}
  \item \prompt{Longer-term improvement and how the client can verify it.}
\end{enumerate}

% SCOPE -----------------------------------------------------------------------
\clearpage
\section{Scope and Approach}
\begin{tabularx}{\textwidth}{@{}lY@{}}
  \toprule
  Client & \ClientName \\
  Testing period & \TestingDates \\
  Team / contact & \TeamName{} / \prompt{Contact} \\
  Starting access & \prompt{External, internal, provided credentials, roles} \\
  \bottomrule
\end{tabularx}

\subsection{Authorized scope}
\begin{tabularx}{\textwidth}{@{}p{0.30\textwidth}Y@{}}
  \toprule
  \textbf{Asset / range} & \textbf{Purpose and boundaries} \\
  \midrule
  \prompt{Hostname, IP, or CIDR} & \prompt{Service and permitted testing} \\
  \prompt{Application or site} & \prompt{URL, roles, and permitted testing} \\
  \bottomrule
\end{tabularx}

\textbf{Excluded activities and assets:}
\prompt{Record exclusions and rules of engagement, including any restrictions
on social engineering, disruption, or accessing sensitive data.}

\subsection{Approach and limitations}
\prompt{Briefly describe reconnaissance, enumeration, manual validation,
and any authorized exploitation. List only activities actually performed.}

\prompt{Record time constraints, unavailable systems, missing credentials,
and untested areas. Explain how each limitation affects the conclusions.}

\subsection{Risk rating}
Findings are prioritized using exploitability and business impact in the
assessed environment. Each finding explains its rating. Informational items
record observations without a demonstrated direct security impact.
\prompt{Adjust this method to the engagement requirements. If using CVSS,
record the version, score, and full vector for each scored finding.}

% FINDINGS OVERVIEW -----------------------------------------------------------
\clearpage
\section{Findings Overview}
\prompt{Summarize the main themes. Order findings by risk. Keep this table
consistent with the detailed findings; the row below is a placeholder.}

% Add one row per finding. longtable continues across pages automatically.
\begin{longtable}{@{}p{0.09\textwidth}p{0.16\textwidth}p{0.38\textwidth}p{0.27\textwidth}@{}}
  \toprule
  \textbf{ID} & \textbf{Severity} & \textbf{Finding} & \textbf{Affected assets} \\
  \midrule
  \endfirsthead
  \toprule
  \textbf{ID} & \textbf{Severity} & \textbf{Finding} & \textbf{Affected assets} \\
  \midrule
  \endhead
  \bottomrule
  \endfoot
  \hyperref[finding:01]{F-01} & \severity{high} &
  \hyperref[finding:01]{[Finding title]} & \prompt{Assets} \\
\end{longtable}

% DETAILED FINDINGS -----------------------------------------------------------
% Copy from \finding through the references prompt to add another finding.
% Assign a unique ID and label, then add a row to the overview above.
% The first finding follows the section title without an extra page break.
\clearpage
\section{Detailed Findings}
\begingroup
\let\clearpage\relax
\finding{F-01}{[Finding title]}{high}{finding:01}
\endgroup

\begin{tabularx}{\textwidth}{@{}lY@{}}
  \toprule
  Affected assets & \prompt{Hosts, IPs, URLs, ports, or accounts} \\
  Preconditions & \prompt{Required access, role, and network position} \\
  Status & \prompt{Open / remediated / partially remediated / not retested} \\
  \bottomrule
\end{tabularx}

\subsubsection*{Description}
\prompt{Explain the weakness, its root cause, and where it occurs. Include
enough context for the client to understand the issue.}

\subsubsection*{Impact and rating rationale}
\prompt{State the demonstrated result and its business impact. Explain why
the likelihood and impact justify the severity. Clearly identify any impact
that is plausible but was not tested.}

\subsubsection*{Reproduction and evidence}
\begin{enumerate}
  \item \prompt{Describe the starting access and exact target.}
  \item \prompt{Provide the request, command, or action used.}
  \item \prompt{Describe the observed result and reference the evidence.}
\end{enumerate}

% lstlisting preserves characters such as _, %, #, and backslashes literally.
\begin{lstlisting}
[Paste a minimal command, HTTP request, or output excerpt here.]
[Redact credentials, session tokens, and unnecessary personal data.]
\end{lstlisting}

\prompt{Evidence ID, capture timestamp with timezone, and what the evidence proves.}

% Optional screenshot: upload an image and uncomment this block.
% \begin{figure}[htbp]
%   \centering
%   \includegraphics[width=\linewidth,height=0.45\textheight,
%                    keepaspectratio]{evidence/f01.png}
%   \caption{F-01: Describe the observed result. Sensitive values are redacted.}
%   \label{fig:f01}
% \end{figure}

\subsubsection*{Remediation}
\prompt{Give a specific corrective action that addresses the root cause.
Include a temporary mitigation if the permanent fix needs more time.}

\subsubsection*{Verification}
\prompt{Explain how to repeat the test after remediation and what secure
behavior should look like. Record retest date and outcome only if retested.}

\subsubsection*{References}
\prompt{Link relevant vendor guidance, CWE, or advisories using
\texttt{\textbackslash url\{https://example.com\}}. Remove if unnecessary.}

% OPTIONAL APPENDIX -----------------------------------------------------------
\clearpage
\appendix
\section{Engagement Notes}
% Remove unused subsections. Keep long raw tool output outside the main report.
\subsection{Attack narrative}
\prompt{If multiple findings formed an attack path, briefly explain the
sequence from initial access to the demonstrated business impact. Reference
finding IDs and distinguish tested steps from hypothetical ones.}

\subsection{Cleanup and handover}
\prompt{List accounts, files, configuration changes, or other artifacts
created during testing. Record removal status and any client action needed.}

\subsection{Supporting evidence}
\prompt{List separately delivered evidence files and their associated finding
IDs. Include timestamps and hashes where useful for traceability.}

\end{document}
